% ======================================================================
%  SECCIONES 5 Y 6 — PARA INTEGRAR EN «tesis primera parte.tex»
%  Insertar entre la sección 4 (Objetivos) y la sección de Referencias.
%
%  REQUIERE añadir al preámbulo (ver instrucciones al final del archivo):
%      \usepackage{amsmath}
%      \usepackage{amssymb}
%      \usepackage{booktabs}
% ======================================================================


% ==================== 5. FUNDAMENTO MATEMÁTICO ====================
\section{Fundamento matemático del problema}

\justifying

Los objetivos planteados exigen medir algo que no es directamente observable:
lo que un modelo de lenguaje representa internamente antes de decidir si
atiende o rechaza una solicitud. Esta sección construye, desde sus elementos
más básicos, el aparato matemático que hace posible esa medición. El orden es
deliberado: primero se establece qué es exactamente el objeto que se observa,
después por qué ese objeto admite una descripción geométrica, y solo entonces
se definen las cantidades que el estudio calcula. Cada expresión formal se
presenta acompañada de la intuición que la motiva y de su lectura en lenguaje
corriente, de modo que el argumento pueda seguirse sin conocimiento previo del
área.

\subsection{Qué se observa: la activación como vector}

Un modelo de lenguaje no procesa texto, procesa números. El texto de entrada se
divide primero en unidades mínimas llamadas \textit{tokens} —fragmentos de
palabra, palabras completas o signos de puntuación— y cada token se convierte
en una lista de números que el modelo puede manipular. A medida que el texto
atraviesa las sucesivas capas de la red, esa lista de números se transforma. En
cualquier punto del recorrido, el estado interno del modelo respecto de un
token determinado es, literalmente, una lista de números reales.

Esa lista recibe el nombre de \textit{activación}. Su longitud es una
característica fija de la arquitectura, denominada dimensión del estado oculto
y denotada $d_{\text{model}}$. En el modelo empleado en este trabajo,
$d_{\text{model}} = 8192$ y la red consta de $L = 80$ capas. Cada token, en
cada capa, tiene asociada una activación de 8192 números.

Conviene detenerse en por qué esta lista se trata como un \textit{vector} y no
simplemente como una colección de mediciones. Un vector no es cualquier lista:
es una lista sobre la cual las operaciones de suma y de multiplicación por un
número tienen un significado consistente. Que las activaciones de un modelo de
lenguaje cumplan esa condición no es una convención adoptada por comodidad
analítica, sino una consecuencia de cómo está construida la arquitectura, y es
lo que se establece a continuación.

\subsection{Por qué el estado interno forma un espacio vectorial}

La arquitectura \textit{transformer} \parencite{vaswani2017} organiza el
cómputo de una manera particular: cada capa no reemplaza el estado que recibe,
sino que le \textit{suma} su contribución. Formalmente, para la capa
$l \in \{1,\dots,L\}$,

\begin{equation}
  x^{(l)} \;=\; x^{(l-1)}
  \;+\; \operatorname{Attn}^{(l)}\!\left(\operatorname{LN}\!\left(x^{(l-1)}\right)\right)
  \;+\; \operatorname{MLP}^{(l)}\!\left(\operatorname{LN}(\cdot)\right),
  \label{eq:residual}
\end{equation}

\noindent
donde $x^{(l)}$ es el estado tras la capa $l$, $\operatorname{Attn}^{(l)}$ es el
bloque de atención de esa capa, $\operatorname{MLP}^{(l)}$ el bloque de
perceptrón multicapa, y $\operatorname{LN}$ una operación de normalización. El
estado inicial $x^{(0)}$ es la representación numérica del texto de entrada.

Leída en palabras, la ecuación~\eqref{eq:residual} dice que el estado tras una
capa es el estado anterior más lo que aporta la atención más lo que aporta el
perceptrón. El elemento decisivo es el signo de suma. Si cada capa sustituyera
el estado recibido, no habría continuidad entre capas y no tendría sentido
hablar de un canal común. Al sumar, en cambio, todas las capas escriben sobre
un mismo soporte compartido, que en la literatura de interpretabilidad se
denomina \textit{flujo residual} \parencite{elhage2021}. Una analogía útil es
la de una grabación multipista: cada músico añade su pista sobre una mezcla
común en lugar de borrar lo que había antes, de modo que la mezcla final
contiene la suma de todos los aportes y es posible preguntarse cuánto hay de
cada instrumento en ella.

Desarrollando la recursión de la ecuación~\eqref{eq:residual} desde la primera
capa hasta la última se obtiene una forma que hace explícita esa acumulación:

\begin{equation}
  x^{(L)} \;=\; \operatorname{Embed}(t)
  \;+\; \sum_{l=1}^{L}\left[\operatorname{Attn}^{(l)} + \operatorname{MLP}^{(l)}\right],
  \label{eq:residual-desarrollado}
\end{equation}

\noindent
donde $t$ denota la secuencia de tokens de entrada. El estado final es la
representación inicial más la suma de todas las contribuciones intermedias.

Esta propiedad aditiva es la que convierte el conjunto de activaciones posibles
en un espacio vectorial genuino de dimensión 8192, dentro del cual la suma de
vectores posee interpretación semántica y las nociones geométricas habituales
—distancia, ángulo, proyección— están bien definidas. Toda la metodología de
este trabajo descansa sobre ella: sin aditividad no sería posible ni proyectar
una activación sobre una dirección ni, como se verá más adelante, intervenir
sumando una dirección al estado interno.

\subsection{Por qué un concepto puede ser una dirección}

Establecido que las activaciones habitan un espacio vectorial, queda por
justificar el supuesto central de la interpretabilidad mecánica: que los
conceptos que el modelo maneja se codifican como \textit{direcciones} de ese
espacio. La formulación es la siguiente: si el modelo representa internamente
una noción determinada, existe un vector $v$ tal que la magnitud del producto
interno $\langle x, v\rangle$ indica en qué grado esa noción está activa en el
estado $x$.

La analogía geográfica resulta esclarecedora. En un mapa, «norte» no es un
lugar sino una dirección; de cualquier ciudad puede preguntarse qué tan al
norte se encuentra y obtener un número. El supuesto sostiene que «peligrosidad»
opera de manera equivalente dentro del modelo: no reside en una neurona ni en
una región concreta, sino que corresponde a una dirección, y de cualquier
activación puede preguntarse cuánto apunta hacia ella.

Este supuesto no es arbitrario, y conviene poder defenderlo porque de él
depende todo lo demás. Su motivación proviene de cómo el modelo \textit{lee} su
propio estado interno para producir la salida. La capa final, denominada
\textit{unembedding}, es una operación lineal:

\begin{equation}
  \operatorname{logits} \;=\; W_U \cdot \operatorname{LN}\!\left(x^{(L)}\right),
  \qquad
  \operatorname{logit}(k) \;=\; \left\langle W_U[k],\; \operatorname{LN}\!\left(x^{(L)}\right)\right\rangle,
  \label{eq:unembedding}
\end{equation}

\noindent
donde $W_U$ es la matriz que proyecta el estado interno sobre el vocabulario y
$W_U[k]$ es la fila correspondiente al token $k$. La ecuación~\eqref{eq:unembedding}
establece que la puntuación asignada a cada token candidato es exactamente un
producto interno entre el estado del modelo y una fila de esa matriz. Dicho de
otro modo, la decisión final se toma mediante productos internos. Resulta
entonces coherente esperar que las nociones que alimentan esa decisión estén
codificadas de una forma que un producto interno pueda leer. El supuesto de
representación lineal queda así motivado por la arquitectura y no impuesto
desde fuera \parencite{elhage2021,zou2023}.

Conviene precisar el alcance de este argumento. Lo que se establece es que la
\textit{lectura} final del estado es lineal, no que el cómputo interno que
genera cada noción lo sea. Se trata de evidencia sólida a favor del supuesto,
no de una demostración, razón por la cual el diseño incorpora una comprobación
empírica específica descrita en la sección~\ref{subsec:validacion}.

\subsection{La dirección de rechazo}
\label{subsec:direccion}

\textcite{arditi2024} mostraron que el rechazo a solicitudes dañinas está
mediado por una sola dirección del espacio de activaciones, y propusieron un
procedimiento para estimarla que este trabajo adopta por dos razones: es
transparente y permite replicación directa de sus resultados.

La idea subyacente es sencilla. Si se dispone de un conjunto de solicitudes
dañinas y otro de solicitudes inofensivas, cada uno produce una nube de puntos
en el espacio de activaciones. Calculando el centro de cada nube y trazando la
flecha que va de la inofensiva a la dañina se obtiene, precisamente, la
dirección que separa ambas condiciones. Formalmente, sean

\begin{equation}
  \mu_H \;=\; \frac{1}{n_H}\sum_{i=1}^{n_H} a_{H,i},
  \qquad
  \mu_B \;=\; \frac{1}{n_B}\sum_{i=1}^{n_B} a_{B,i}
  \label{eq:centroides}
\end{equation}

\noindent
los centroides de las activaciones dañinas e inofensivas respectivamente,
donde $a_{H,i}$ y $a_{B,i}$ denotan activaciones individuales y $n_H$, $n_B$
el número de observaciones en cada grupo. La dirección de rechazo se define
como la diferencia de esos centroides, normalizada a longitud unitaria:

\begin{equation}
  d_{\text{chat}} \;=\; \frac{\mu_H - \mu_B}{\lVert \mu_H - \mu_B \rVert_2}.
  \label{eq:direccion}
\end{equation}

El símbolo $\lVert \cdot \rVert_2$ denota la longitud euclídea del vector, es
decir, la raíz cuadrada de la suma de los cuadrados de sus componentes. Dividir
por ella garantiza que $d_{\text{chat}}$ tenga longitud exactamente uno, lo que
lo convierte en una regla de medir de escala fija y permite comparar ángulos
con otras direcciones más adelante.

Este estimador, conocido como diferencia de medias, no es una simplificación
adoptada por conveniencia. El discriminante lineal óptimo para separar dos
clases, formulado por \textcite{fisher1936}, es

\begin{equation}
  w^{*} \;=\; \Sigma^{-1}\left(\mu_H - \mu_B\right),
  \label{eq:fisher}
\end{equation}

\noindent
donde $\Sigma$ es la matriz de covarianza intraclase, que describe cómo varían
conjuntamente las distintas dimensiones dentro de cada grupo. El factor
$\Sigma^{-1}$ cumple la función de descontar aquellas direcciones que son
intrínsecamente ruidosas antes de efectuar la comparación. Si todas las
direcciones fuesen igualmente ruidosas —condición conocida como covarianza
isotrópica, $\Sigma \approx \sigma^2 I$— no habría nada que descontar y la
ecuación~\eqref{eq:fisher} se reduciría a $w^{*} \propto \mu_H - \mu_B$, es
decir, exactamente a la ecuación~\eqref{eq:direccion}.

Adoptar ese supuesto no responde a una preferencia teórica sino a una
imposibilidad práctica. Estimar $\Sigma$ requeriría del orden de
$d_{\text{model}}^{2}/2 \approx 3{,}4 \times 10^{7}$ parámetros, mientras que
el estudio dispone de 42 observaciones por condición. La matriz de covarianza
muestral resultante es singular, con rango a lo sumo $n-1 = 41$, y por tanto no
admite inversión. En este régimen la diferencia de medias es el estimador
disponible que está bien condicionado, y la limitación que ello impone se
retoma en la sección~\ref{subsec:supuestos}.

\subsection{La proyección: de ocho mil dimensiones a un número}

Disponer de una dirección permite reducir cada activación a una sola cantidad
interpretable. Para una activación $a$ y una dirección unitaria $d$, se define
la proyección como

\begin{equation}
  P(a) \;=\; \langle a, d\rangle \;=\; \sum_{j=1}^{d_{\text{model}}} a_j\, d_j,
  \label{eq:proyeccion}
\end{equation}

\noindent
donde el subíndice $j$ recorre las 8192 componentes de ambos vectores. En
palabras: se multiplican componente a componente y se suman todos los
productos.

Geométricamente, la proyección corresponde a la sombra del vector $a$ sobre el
eje que define $d$. Si se ilumina perpendicularmente a ese eje, la sombra
proyectada tiene una longitud determinada, y esa longitud es $P(a)$. El valor
lleva signo: resulta negativo cuando el vector apunta en sentido opuesto al de
la dirección. Que $d$ tenga longitud unitaria es lo que permite esta lectura
directa; sin normalizar, la escala de la sombra sería arbitraria y las
comparaciones perderían sentido.

La ecuación~\eqref{eq:proyeccion} realiza una compresión considerable: convierte
un objeto de 8192 dimensiones en un único número real, conservando íntegramente
la variable de interés. La interpretación operativa que la sustenta es que un
valor alto de $P(a)$ indica que el mecanismo de rechazo está activo y el modelo
tenderá a negarse, mientras que un valor bajo indica lo contrario. Todo el
análisis estadístico posterior opera sobre este escalar.

\subsection{La brecha de proyección}
\label{subsec:deltap}

Con las piezas anteriores puede formularse ya la cantidad que responde al
tercer objetivo específico. Si una misma solicitud dañina se presenta al modelo
en formato conversacional y en formato agéntico, y en cada caso se mide cuánto
se activa la dirección de rechazo, la diferencia entre ambas mediciones
cuantifica el efecto del formato. Esa diferencia se denomina brecha de
proyección y se define como

\begin{equation}
  \Delta P
  \;=\;
  \frac{1}{n}\sum_{i=1}^{n}\left\langle a_{\text{chat},i}, d_{\text{chat}}\right\rangle
  \;-\;
  \frac{1}{n}\sum_{i=1}^{n}\left\langle a_{\text{agent},i}, d_{\text{chat}}\right\rangle,
  \label{eq:deltap}
\end{equation}

\noindent
donde $a_{\text{chat},i}$ y $a_{\text{agent},i}$ son las activaciones
correspondientes a la solicitud $i$ en cada formato, y $n$ es el número de
solicitudes.

La expresión admite una reescritura que ilumina su significado. Puesto que el
producto interno es una operación lineal, promediar las proyecciones equivale a
proyectar el promedio, de donde

\begin{equation}
  \Delta P
  \;=\;
  \left\langle\; \bar{a}_{\text{chat}} - \bar{a}_{\text{agent}},\;\; d_{\text{chat}} \;\right\rangle,
  \label{eq:deltap-lineal}
\end{equation}

\noindent
siendo $\bar{a}$ el centroide de la condición respectiva. La
ecuación~\eqref{eq:deltap-lineal} muestra que $\Delta P$ es la proyección, sobre
el eje de rechazo, del vector que describe cómo se desplaza la nube de
activaciones al cambiar de formato. El cambio de formato mueve el estado
interno en alguna dirección del espacio; $\Delta P$ mide qué porción de ese
movimiento ocurre precisamente a lo largo del eje que gobierna la decisión de
rechazar. Un desplazamiento perpendicular a $d_{\text{chat}}$, por grande que
sea, no contribuye a la brecha.

La lectura de la cantidad es directa. Un valor positivo y estadísticamente
significativo indica que el formato agéntico desplaza las activaciones en
sentido contrario a la activación del rechazo, lo que constituiría evidencia de
que la defensa se atenúa. Un valor próximo a cero indicaría que la
representación se mantiene y que el fallo conductual documentado por
\textcite{andriushchenko2024} se origina en etapas posteriores del
procesamiento. Ambos desenlaces son informativos.

\subsection{Cuando la dirección deja de servir: la rotación}
\label{subsec:rotacion}

Una brecha positiva admite al menos dos explicaciones que conviene distinguir,
porque tienen implicaciones distintas para el quinto objetivo específico. Puede
ocurrir que el modelo deje de codificar la peligrosidad de la solicitud, o
puede ocurrir que siga codificándola pero en una dirección distinta, que el
mecanismo original ya no consulta. Separar ambas posibilidades requiere una
segunda medición.

Para ello se estima una dirección adicional, denotada $w_{\text{agent}}$,
entrenando un clasificador lineal exclusivamente sobre activaciones de formato
agéntico. El clasificador empleado es una regresión logística, que modela la
probabilidad de que una activación corresponda a contenido dañino mediante

\begin{equation}
  \Pr(y = 1 \mid z) \;=\; \sigma\!\left(w^{\!\top} z + b\right),
  \qquad
  \sigma(u) \;=\; \frac{1}{1 + e^{-u}},
  \label{eq:logistica}
\end{equation}

\noindent
donde $y$ vale uno para las solicitudes dañinas y cero para las inofensivas,
$z$ es la activación, $w$ el vector de pesos, $b$ un término independiente y
$\sigma$ la función logística, que transforma cualquier número real en un valor
entre cero y uno interpretable como probabilidad. El vector $w$ resultante es
perpendicular a la frontera que el clasificador traza entre ambas clases y,
normalizado, constituye la dirección que mejor discrimina la peligrosidad en
contexto agéntico.

Comparar ambas direcciones equivale a medir el ángulo que las separa. Siendo
las dos de longitud unitaria, el coseno de ese ángulo es simplemente su
producto interno:

\begin{equation}
  \cos\theta \;=\; \left\langle w_{\text{agent}},\, d_{\text{chat}}\right\rangle,
  \qquad
  \theta \;=\; \arccos\!\left(\cos\theta\right).
  \label{eq:coseno}
\end{equation}

Un coseno próximo a uno indicaría que ambas direcciones coinciden y que no hay
reorganización alguna; un coseno próximo a cero indicaría que son
perpendiculares, es decir, que la información sobre peligrosidad se ha
reubicado por completo en un eje que el mecanismo original no examina; y
valores intermedios corresponderían a una rotación parcial.

Interpretar esta cantidad exige una advertencia que resulta contraintuitiva y
que es propia de los espacios de dimensión elevada. En $\mathbb{R}^{8192}$, dos
vectores tomados al azar son casi perpendiculares entre sí con altísima
probabilidad: su coseno se distribuye aproximadamente como

\begin{equation}
  \cos\theta_{\text{azar}} \;\sim\; \mathcal{N}\!\left(0,\; \tfrac{1}{d_{\text{model}}}\right),
  \qquad
  \sigma_{\text{azar}} \;=\; \frac{1}{\sqrt{d_{\text{model}}}} \;\approx\; 0{,}011.
  \label{eq:coseno-azar}
\end{equation}

La ecuación~\eqref{eq:coseno-azar} proporciona la escala de referencia
indispensable para leer cualquier coseno observado. Sin ella, un valor de
$0{,}2$ parecería indicar ausencia de alineación; con ella se comprende que
equivale a unas dieciocho desviaciones estándar por encima de lo esperable por
azar. En dimensión elevada, la ausencia de alineación no se sitúa
intuitivamente cerca de cero: se sitúa a una distancia de cero que hay que
calcular.

Finalmente, es necesario distinguir dos propiedades que suelen confundirse y
cuya separación constituye un aspecto central de este trabajo. Una dirección
puede conservar intacta su capacidad de \textit{ordenar} correctamente las
solicitudes —situando siempre las dañinas por encima de las inofensivas— y
simultáneamente perder el \textit{nivel absoluto} de activación necesario para
que el mecanismo de rechazo se dispare. La primera propiedad se mide mediante
el área bajo la curva característica del receptor y la segunda mediante
$\Delta P$. Que ambas puedan disociarse implica que el modelo podría seguir
distinguiendo perfectamente lo dañino de lo inofensivo mientras deja de actuar
en consecuencia, escenario que este estudio está diseñado para detectar.


% ==================== 6. DISEÑO METODOLÓGICO ====================
\section{Diseño metodológico}

\justifying

\subsection{Enfoque, tipo y alcance}

La investigación adopta un enfoque cuantitativo, en tanto la totalidad de sus
inferencias descansa sobre magnitudes numéricas medidas en el espacio de
activaciones y contrastadas mediante pruebas estadísticas formales. El tipo de
estudio es experimental: la variable independiente se manipula deliberadamente,
manteniendo constantes el contenido textual de la solicitud, los pesos del
modelo y los parámetros de generación. No se trata de un estudio observacional,
pues las condiciones comparadas no se registran tal como ocurren en la
naturaleza sino que se construyen expresamente.

El alcance es explicativo. El trabajo no se limita a describir la brecha
conductual entre formatos, ya documentada \parencite{andriushchenko2024}, sino
que busca identificar el mecanismo representacional que la acompaña. Dos
elementos sostienen esa pretensión: el control experimental del contenido
mediante emparejamiento, descrito enseguida, y la incorporación de un
procedimiento de intervención sobre las activaciones orientado a contrastar
suficiencia causal y no únicamente asociación. El estudio es además de corte
transversal, pues todas las mediciones se obtienen sobre un mismo estado de los
pesos del modelo, sin reentrenamiento ni seguimiento temporal.

\subsection{El diseño pareado y por qué es la decisión central}

El primer objetivo específico exige aislar el efecto del formato manteniendo
constante el contenido de la instrucción. La forma de lograrlo es un diseño
intrasujeto de medidas repetidas, en el cual cada solicitud se presenta al
modelo bajo todas las condiciones experimentales en lugar de asignar
solicitudes distintas a condiciones distintas.

La razón es la misma por la que, para evaluar un tratamiento médico, comparar a
un paciente consigo mismo antes y después resulta más informativo que comparar
dos pacientes diferentes: en el segundo caso las diferencias individuales
contaminan la medición. Aquí el papel del paciente lo desempeña el texto de la
solicitud. Denotando por $a_{\text{chat},i}$ y $a_{\text{agent},i}$ las
activaciones de la solicitud $i$ en cada formato, y siendo el mensaje del
usuario idéntico en ambos casos, se cumple

\begin{equation}
  a_{\text{agent},i} - a_{\text{chat},i}
  \;=\;
  f(\text{formato}) \;+\; \underbrace{0}_{\text{aporte del contenido}},
  \label{eq:cancelacion}
\end{equation}

\noindent
donde $f(\cdot)$ representa el efecto atribuible al andamiaje del formato. La
ecuación~\eqref{eq:cancelacion} expresa que, al restar dentro del par, el
contenido se cancela algebraicamente y toda diferencia observada queda
imputada al formato. El contenido deja de ser una variable capaz de confundir
la comparación y pasa a ser una constante del diseño.

El emparejamiento aporta además una ganancia estadística que conviene explicitar.
Para dos variables aleatorias cualesquiera,

\begin{equation}
  \operatorname{Var}(A - B)
  \;=\;
  \operatorname{Var}(A) + \operatorname{Var}(B) - 2\operatorname{Cov}(A,B),
  \label{eq:varianza-diferencia}
\end{equation}

\noindent
donde $\operatorname{Cov}(A,B)$ mide en qué grado ambas cantidades varían de
forma conjunta. Como las dos condiciones comparten el texto de la solicitud,
esa covarianza es elevada y, al sustraerse, reduce la varianza del contraste.
En términos prácticos, el diseño pareado alcanza la misma potencia estadística
con un número menor de observaciones, lo que resulta determinante dado el costo
computacional de cada medición.

\subsection{Corpus contrastivo}

El corpus fue construido específicamente para este estudio, dado que los
conjuntos de datos disponibles no satisfacen el requisito de emparejamiento:
en el \textit{benchmark} AgentHarm, por ejemplo, las solicitudes empleadas en
cada condición son textos distintos, de modo que cualquier diferencia observada
resultaría indistinguible de una diferencia de contenido
\parencite{andriushchenko2024}. Su elaboración se realizó en el marco de una
\textit{hackathon} de seguridad en inteligencia artificial.

\medskip
\noindent
\textbf{[Pendiente por completar: responsables de la construcción del corpus,
protocolo de redacción de las solicitudes, criterios operativos aplicados para
asignar el grado de explicitud, número de evaluadores que participaron en el
etiquetado y coeficiente de concordancia obtenido entre ellos.]}
\medskip

El corpus consta de 42 solicitudes base, cada una presente en cuatro variantes
resultantes de cruzar formato (conversacional o agéntico) con contenido
(dañino o inofensivo), lo que arroja un total de 168 registros. Las solicitudes
se distribuyen en siete categorías de daño —fraude, cibercrimen,
desinformación, acoso, odio, derechos de autor y drogas— y en tres grados de
explicitud de la intención dañina, definidos operativamente así: explícito,
cuando la intención es manifiesta en el léxico de la solicitud; contextual,
cuando el daño se infiere del contexto y no de los términos literales; y
encubierto, cuando la acción dañina se presenta bajo un marco de legitimidad
profesional. El cuadro~\ref{tab:corpus} resume la distribución.

\begin{table}[htbp]
  \centering
  \small
  \caption{Distribución del corpus, expresada en solicitudes base.}
  \label{tab:corpus}
  \begin{tabular}{@{}lr@{\hspace{3em}}lr@{}}
    \toprule
    \textbf{Categoría} & \textbf{$n$} & \textbf{Explicitud} & \textbf{$n$} \\
    \midrule
    Fraude            & 7 & Explícito   & 18 \\
    Cibercrimen       & 7 & Contextual  & 14 \\
    Desinformación    & 6 & Encubierto  & 10 \\
    Acoso             & 6 &             &    \\
    Odio              & 6 &             &    \\
    Derechos de autor & 5 &             &    \\
    Drogas            & 5 &             &    \\
    \midrule
    \textbf{Total}    & \textbf{42} & \textbf{Total} & \textbf{42} \\
    \bottomrule
  \end{tabular}
\end{table}

Dado que la validez de la ecuación~\eqref{eq:cancelacion} depende íntegramente
de que el texto sea idéntico entre condiciones, se ejecutó una verificación
programática que comparó carácter a carácter el mensaje del usuario en ambos
formatos para cada combinación de solicitud base y contenido. La verificación
confirmó la identidad textual en la totalidad de los 84 pares, lo que valida la
premisa sobre la que se apoya el diseño.

La condición inofensiva en formato agéntico, que podría parecer redundante,
cumple una función de control indispensable. Si el formato agéntico desplazara
todas las activaciones por igual —por efecto, por ejemplo, de una mayor
longitud del contexto— ese desplazamiento carecería de relación con la
seguridad. Solo disponiendo del par inofensivo puede distinguirse un efecto
específico sobre la representación de peligrosidad de un corrimiento global
inducido por el formato.

\subsection{Operacionalización del formato agéntico}

La distinción entre ambos formatos se implementó mediante la instrucción de
sistema que precede al mensaje del usuario. En la condición conversacional
dicha instrucción describe al modelo como un asistente que responde
solicitudes. En la condición agéntica lo describe como un asistente con acceso
a herramientas, al que se indica que responda mediante una invocación de
herramienta cuando corresponda.

\medskip
\noindent
\textbf{[Observación que debe consignarse con precisión: en la ejecución
realizada, las especificaciones de las herramientas asociadas a cada solicitud
—almacenadas en el corpus como estructuras independientes— no llegaron a
incorporarse al texto del \textit{prompt}. Una verificación programática
confirmó que en ninguno de los 84 registros agénticos la instrucción de sistema
menciona las herramientas correspondientes. En consecuencia, lo efectivamente
manipulado en dicha ejecución fue el encuadre de rol y la instrucción de
invocar herramientas, no la presencia de las especificaciones en el contexto.
Esta precisión no invalida la comparación —la manipulación resulta, de hecho,
más económica, pues descarta que un eventual efecto se deba al volumen de texto
añadido— pero sí acota lo que puede afirmarse: la evidencia obtenida concierne
al encuadre de rol. Determinar la contribución específica de las
especificaciones de herramientas requiere una ejecución adicional en la que
estas se inyecten efectivamente, prevista como trabajo inmediato.]}
\medskip

\subsection{Modelo y extracción de las activaciones}

El sistema analizado es un modelo de lenguaje autorregresivo con ajuste por
instrucciones, de 80 capas y dimensión de estado oculto 8192. Su elección
respondió a tres criterios: disponer de un entrenamiento de seguridad
suficientemente robusto para que el mecanismo de rechazo resulte observable,
contar con soporte nativo de invocación de herramientas, y tratarse de un
modelo ampliamente estudiado, lo que facilita la comparación con la literatura.

La recolección de activaciones se realizó mediante \textit{ganchos de
propagación}, funciones que el entorno de ejecución invoca automáticamente tras
el cómputo de cada capa y que permiten registrar su salida sin alterar el
resultado del modelo. Su papel durante esta etapa es estrictamente pasivo:
observan sin modificar, del mismo modo que un caudalímetro mide sin perturbar
el flujo de la tubería.

De la secuencia completa se retiene una única posición: la del último token.
La justificación es estructural. La atención del modelo es causal, condición
que se implementa mediante una máscara

\begin{equation}
  M_{ij} \;=\;
  \begin{cases}
    0        & \text{si } j \le i,\\[2pt]
    -\infty  & \text{si } j > i,
  \end{cases}
  \label{eq:mascara-causal}
\end{equation}

\noindent
donde $i$ indexa la posición que consulta y $j$ la posición consultada. El valor
$-\infty$ se suma a las puntuaciones antes de normalizarlas, de modo que el peso
asignado a cualquier posición futura resulta exactamente nulo. Se sigue de la
ecuación~\eqref{eq:mascara-causal} que cada token únicamente integra información
de los que le preceden, y por tanto que el último token es el único que ha
incorporado la totalidad del contexto. Es además el estado desde el cual se
computa la distribución del siguiente token según la
ecuación~\eqref{eq:unembedding}, esto es, el punto exacto en que la decisión se
produce. Una consideración adicional refuerza la elección: al variar la longitud
de las solicitudes, la última posición es la única cuyo significado es
equivalente entre condiciones.

Las activaciones se convirtieron a precisión de 32 bits antes del análisis. La
razón es que la ecuación~\eqref{eq:proyeccion} acumula 8192 sumandos, y en
precisión reducida el error de redondeo acumulado podría alcanzar el orden de
magnitud de la señal que se pretende medir.

\subsection{Variables}

El cuadro~\ref{tab:variables} distingue explícitamente entre lo que se
manipula, lo que se mide, lo que se mantiene fijo y lo que se obtiene como
resultado del procedimiento.

\begin{table}[htbp]
  \centering
  \small
  \caption{Clasificación de las variables del estudio.}
  \label{tab:variables}
  \begin{tabular}{@{}llp{5.6cm}@{}}
    \toprule
    \textbf{Clase} & \textbf{Elemento} & \textbf{Naturaleza} \\
    \midrule
    Independiente & Formato de presentación & Categórica: conversacional / agéntico \\
    Estratificación & Grado de explicitud & Categórica ordinal: explícito / contextual / encubierto \\
    \addlinespace
    Dependiente & Proyección $P(a)$ & Continua, en unidades del espacio de activación \\
    Dependiente & Brecha $\Delta P$ & Continua, mismas unidades \\
    Dependiente & Alineación $\cos\theta$ & Adimensional, entre $-1$ y $1$ \\
    \addlinespace
    Controlada & Texto de la solicitud & Idéntico entre formatos \\
    Controlada & Pesos del modelo & Constantes; no se reentrena \\
    Controlada & Posición tokénica & Fija: último token \\
    Controlada & Generación & Determinista, sin muestreo \\
    \addlinespace
    Resultado & $d_{\text{chat}}$, $w_{\text{agent}}$ & Vectores unitarios en $\mathbb{R}^{8192}$ \\
    \bottomrule
  \end{tabular}
\end{table}

Una precisión sobre las unidades resulta necesaria para interpretar
correctamente los valores. La proyección se expresa en las unidades arbitrarias
del espacio de activación, y la norma del estado interno crece
sistemáticamente con la profundidad de la capa. En consecuencia, los valores de
$P$ y $\Delta P$ \textit{no son directamente comparables entre capas} sin una
normalización previa. Las cantidades adimensionales —el coseno y el área bajo
la curva— sí admiten esa comparación, y son las que se emplean cuando el
interés recae sobre la evolución del fenómeno a lo largo de la red.

\subsection{Inferencia estadística}

El contraste de la hipótesis se apoya en tres pruebas que descansan sobre
supuestos distintos, de modo que la coincidencia entre ellas ofrece garantías
que ninguna proporcionaría por separado.

La primera es la prueba $t$ para muestras pareadas. Definiendo la diferencia
dentro de cada par como $D_i = P_{\text{chat},i} - P_{\text{agent},i}$, el
estadístico de contraste es

\begin{equation}
  t \;=\; \frac{\bar{D}}{s_D/\sqrt{n}},
  \qquad \text{con } n-1 \text{ grados de libertad},
  \label{eq:t-pareado}
\end{equation}

\noindent
donde $\bar{D}$ es la media de las diferencias y $s_D$ su desviación estándar.
El cociente compara la magnitud del cambio promedio con la variabilidad de ese
cambio: un valor elevado indica que la caída es consistente entre solicitudes y
no producto del azar. Conviene señalar que esta prueba supone normalidad de las
\textit{diferencias}, no de las proyecciones, supuesto considerablemente más
débil.

La segunda prueba prescinde por completo de supuestos distribucionales. Bajo la
hipótesis nula de que el formato no influye, las etiquetas de condición serían
intercambiables dentro de cada par; puede entonces construirse la distribución
nula invirtiendo aleatoriamente el signo de cada diferencia y repitiendo el
procedimiento un número elevado de veces. El valor de probabilidad empírico es

\begin{equation}
  \hat{p} \;=\; \frac{1}{B}\sum_{b=1}^{B}
  \mathbb{1}\!\left[\Delta P^{(b)}_{\text{nula}} \;\ge\; \Delta P_{\text{obs}}\right],
  \label{eq:permutacion}
\end{equation}

\noindent
donde $B$ es el número de permutaciones, $\mathbb{1}[\cdot]$ vale uno cuando la
condición se cumple y cero en caso contrario, y $\Delta P_{\text{obs}}$ es la
brecha efectivamente observada. Se emplea $B = 10\,000$, lo que fija la
resolución del procedimiento: un valor reportado como nulo debe leerse como
inferior a $1/B$, nunca como cero exacto.

La tercera es un intervalo de confianza por remuestreo \parencite{efron1979},
que aporta lo que ninguna prueba de hipótesis proporciona: una estimación de la
precisión. Se extraen repetidamente muestras con reemplazo de los pares
observados y se recalcula la brecha en cada réplica, tomando después los
percentiles correspondientes al nivel de confianza deseado. Resulta esencial
que el mismo vector de índices se aplique a ambas condiciones, pues extraer
índices independientes destruiría la estructura de emparejamiento y estimaría
una cantidad distinta de la buscada.

Junto a las pruebas se reportan tamaños de efecto, que informan sobre la
magnitud del fenómeno con independencia del tamaño muestral \parencite{cohen1988}.
Se emplean dos versiones que no son intercambiables y cuya distinción debe
explicitarse: la pareada, $\bar{D}/s_D$, que responde a qué tan consistente es
el cambio dentro de cada par; y la agrupada, que trata ambas condiciones como
independientes y responde a qué tan separadas están las dos distribuciones. Ambas
se relacionan mediante la correlación intra-par $r$ según

\begin{equation}
  d_{\text{pareado}} \;=\; \frac{d_{\text{agrupado}}}{\sqrt{2\,(1-r)}},
  \label{eq:relacion-d}
\end{equation}

\noindent
expresión que permite recuperar $r$ a partir de ambos valores y constituye, por
tanto, un diagnóstico directo de en qué medida el emparejamiento cumplió su
función.

Finalmente, la capacidad de una dirección para ordenar correctamente
solicitudes dañinas e inofensivas se evalúa mediante el área bajo la curva
característica del receptor, cuya interpretación probabilística es exacta:
corresponde a la probabilidad de que una solicitud dañina tomada al azar reciba
una proyección mayor que una inofensiva tomada al azar. Un valor de $0{,}5$
equivale a decidir al azar y uno de $1{,}0$ a separación perfecta.

El cuarto objetivo específico se atiende replicando este conjunto de
procedimientos de forma independiente dentro de cada grado de explicitud, lo
que permite determinar si el efecto depende de cuán manifiesta sea la intención
dañina.

\subsection{Procedimientos de validación}
\label{subsec:validacion}

Se contemplan tres estrategias de validación, cada una dirigida a una pregunta
distinta.

La validación en el espacio de activaciones examina si la dirección estimada
constituye un eje real de variación o un artefacto del ruido muestral. Para
ello se contrasta $d_{\text{chat}}$ con la primera componente principal de los
datos y se comparan además dos clasificadores: uno lineal y uno no lineal. Este
segundo contraste es la comprobación empírica del supuesto de representación
lineal anunciada anteriormente: si ambos alcanzan un desempeño equivalente, la
información se organiza efectivamente como una dirección; si el no lineal
resulta marcadamente superior, la información existe pero no admite descripción
mediante un vector único, lo que obligaría a revisar el marco formal adoptado.

La validación causal responde a una objeción de fondo: todo lo anterior es
observacional, y la asociación entre proyección y comportamiento no demuestra
que la primera cause el segundo. Para contrastarlo se implementa una
intervención en tiempo de inferencia consistente en sumar la dirección al flujo
residual durante la generación,

\begin{equation}
  x^{(l)} \;\longleftarrow\; x^{(l)} \;+\; \alpha\, d,
  \label{eq:intervencion}
\end{equation}

\noindent
donde $\alpha$ regula la intensidad de la perturbación \parencite{turner2023,zou2023}.
Esta operación es posible precisamente por la aditividad establecida en la
ecuación~\eqref{eq:residual}: sumar un vector al flujo residual en una capa hace
que ese aporte se propague a las capas siguientes. El criterio de contraste es
la existencia de un valor mínimo de $\alpha$ que restituya el rechazo; si
existe, la dirección es causalmente suficiente para inducirlo. Debe precisarse
que este procedimiento demuestra suficiencia, no necesidad: no descarta que
existan otras vías capaces de producir el mismo efecto.

La validación conductual, por último, contrasta sobre un corpus complementario
si los patrones de activación se corresponden con el comportamiento
efectivamente emitido por el modelo.

El quinto objetivo específico se aborda evaluando la viabilidad de un
mecanismo de monitoreo que opere sobre la dirección estimada. Su formulación es
deliberadamente simple: se calcula la proyección de la activación y se emite
una alerta cuando supera un umbral calibrado. Su interés práctico radica en el
costo: la operación requiere $d_{\text{model}}$ multiplicaciones, cantidad
constante respecto del número de herramientas de que disponga el agente, lo que
contrasta con los esquemas de filtrado que exigen una regla por herramienta y
cuyo costo crece con la superficie de exposición.

\subsection{Estimación con corrección de sesgo}

La combinación de un tamaño muestral reducido con una dimensionalidad muy
elevada introduce sesgos que no son evidentes y que, de no corregirse, conducen
a sobrestimar sistemáticamente los efectos. El procedimiento incorpora tres
correcciones cuya necesidad se verificó mediante simulación con datos de verdad
conocida.

La primera atiende una circularidad. La dirección $d_{\text{chat}}$ se estima a
partir de las activaciones conversacionales dañinas, y esas mismas activaciones
se proyectan después sobre ella para calcular el primer término de la
ecuación~\eqref{eq:deltap}. Como el centroide muestral incorpora una componente
de ruido que se transfiere a la dirección, proyectarlo de vuelta recupera parte
de ese ruido e infla la medición; la condición agéntica, que no participa en la
estimación, no experimenta esa inflación, de modo que la brecha absorbe la
diferencia. La corrección consiste en estimación cruzada
\parencite{chernozhukov2018}: se divide la muestra en subconjuntos, se estima
la dirección excluyendo uno de ellos y se proyecta únicamente ese subconjunto
excluido, garantizando que ninguna proyección se realice sobre una dirección
ajustada con la misma observación.

La segunda atiende la evaluación de $w_{\text{agent}}$ y del monitor. Cuando el
número de observaciones es inferior a la dimensión del espacio, el teorema de
\textcite{cover1965} garantiza que cualquier asignación de etiquetas es
separable linealmente. Con 84 observaciones en 8192 dimensiones, un desempeño
perfecto sobre los mismos datos empleados para el ajuste es una certeza
geométrica y no constituye evidencia alguna sobre la información contenida en
ellos. La corrección consiste en evaluar el clasificador y calibrar el umbral
exclusivamente sobre observaciones que no intervinieron en el ajuste.

La tercera atiende la medición del ángulo. La estimación de $w_{\text{agent}}$
a partir de una muestra reducida está dominada por el ruido, y en dimensión
elevada ese ruido es prácticamente perpendicular a cualquier dirección fija, lo
que atenúa el coseno hacia cero. El ángulo medido resulta así una cota superior
del verdadero. Para cuantificar esta limitación se estima la reproducibilidad
de $w_{\text{agent}}$ comparando estimaciones obtenidas de mitades disjuntas de
la muestra, cantidad que determina con qué precisión puede reportarse el ángulo,
o si puede reportarse en absoluto \parencite{spearman1904}.

\subsection{Supuestos y limitaciones}
\label{subsec:supuestos}

El procedimiento descansa sobre un conjunto de supuestos cuyo incumplimiento
comprometería la validez de las inferencias, y que conviene enunciar sin
atenuantes. Se asume que los conceptos relevantes se codifican como direcciones
del espacio de activaciones y que el rechazo está mediado por una dirección
única y no por un subespacio de dimensión mayor. Se asume la isotropía de la
covarianza intraclase, necesaria para que la ecuación~\eqref{eq:direccion}
aproxime a la ecuación~\eqref{eq:fisher}. Se asume que el estado en la última
posición contiene la información determinante de la decisión. Y se asume
intercambiabilidad dentro de cada par, supuesto sobre el que se construye la
distribución nula de la ecuación~\eqref{eq:permutacion}.

Las limitaciones del diseño son las siguientes. El corpus es de construcción
intencional y no procede de un muestreo probabilístico, por lo que las
inferencias conciernen al efecto del formato dentro del conjunto de estímulos
empleado y su extensión a la población de solicitudes posibles constituye un
argumento de cobertura, no de representatividad estadística. El tamaño muestral,
42 pares en un espacio de 8192 dimensiones, sitúa toda estimación vectorial en
el régimen de alta dimensionalidad con las consecuencias descritas en el
apartado anterior. El análisis abarcó un subconjunto de capas que cubre el
tercio inicial de la red, lo que impide formular afirmaciones sobre el
comportamiento en capas profundas. El estudio se aplica a un modelo principal,
de modo que la dependencia de los hallazgos respecto de la escala o de la
familia arquitectónica no puede establecerse con esta evidencia. El corpus está
redactado en un solo idioma. Y, conforme a lo consignado anteriormente, la
ejecución realizada manipuló el encuadre de rol sin incorporar las
especificaciones de herramientas al contexto, lo que acota el alcance de las
afirmaciones sobre la condición agéntica hasta que se complete la ejecución
prevista.

Finalmente, el diseño aborda el daño contenido semánticamente en una única
solicitud. Quedan explícitamente fuera de alcance el daño que emerge de la
historia conversacional acumulada y el que se distribuye entre sesiones
distintas, casos en los que la información necesaria no está disponible en una
sola pasada del modelo y que requieren aproximaciones de otra naturaleza.

\subsection{Consideraciones éticas}

La investigación no involucra participantes humanos ni datos personales, por lo
que no requiere consentimiento informado ni tratamiento de información
identificable. El corpus contiene solicitudes formuladas deliberadamente con
contenido dañino, cuya única finalidad es activar los mecanismos de rechazo del
modelo en condiciones controladas. Dichas solicitudes no se ejecutan contra
sistemas reales: las herramientas son especificaciones declarativas sin
implementación funcional, y el procedimiento analiza representaciones internas
sin materializar acción alguna.

La finalidad del trabajo es defensiva, en tanto busca caracterizar una
vulnerabilidad representacional para hacer posible su mitigación. Se reconoce,
no obstante, el carácter de doble uso del conocimiento generado, pues la
descripción del mecanismo que atenúa el rechazo podría orientar su elusión
deliberada. Dos consideraciones atenúan esta tensión: el fenómeno se encontraba
ya documentado públicamente en el plano conductual con anterioridad a este
trabajo \parencite{andriushchenko2024}, y la contribución se orienta a la
detección antes que a la evasión.

\medskip
\noindent
\textbf{[Pendiente por completar: constancia de aprobación del comité de ética
institucional, si el programa lo requiere; condiciones de licencia y política de
uso aceptable del modelo empleado; y política de publicación del corpus, con la
justificación correspondiente.]}


% ======================================================================
%  INSTRUCCIONES DE INTEGRACIÓN
% ----------------------------------------------------------------------
%
%  1) AÑADIR AL PREÁMBULO, después de \usepackage{lmodern}:
%
%       % ==================== MATEMÁTICAS Y TABLAS ====================
%       \usepackage{amsmath}
%       \usepackage{amssymb}
%       \usepackage{booktabs}
%
%     amsmath  -> entornos equation, cases, align y \text dentro de fórmulas
%     amssymb  -> símbolos \mathbb (usado en R^8192 y en la indicadora)
%     booktabs -> reglas \toprule, \midrule, \bottomrule de los dos cuadros
%
%  2) INSERTAR el contenido de este archivo entre la sección 4 (Objetivos)
%     y el bloque «% ==================== REFERENCIAS ====================».
%
%  3) AÑADIR AL bibliography las ocho entradas nuevas que aparecen abajo,
%     manteniendo el orden alfabético del listado existente.
%
% ======================================================================
%
%  ENTRADAS BIBLIOGRÁFICAS NUEVAS
%  (copiar dentro del entorno thebibliography ya existente)
%
% \bibitem[Chernozhukov et al.(2018)]{chernozhukov2018}
% Chernozhukov, V., Chetverikov, D., Demirer, M., Duflo, E., Hansen, C.,
% Newey, W., \& Robins, J. (2018). Double/debiased machine learning for
% treatment and structural parameters. \textit{The Econometrics Journal},
% 21(1), C1--C68.
%
% \bibitem[Cohen(1988)]{cohen1988}
% Cohen, J. (1988). \textit{Statistical power analysis for the behavioral
% sciences} (2.ª ed.). Lawrence Erlbaum Associates.
%
% \bibitem[Cover(1965)]{cover1965}
% Cover, T. M. (1965). Geometrical and statistical properties of systems of
% linear inequalities with applications in pattern recognition.
% \textit{IEEE Transactions on Electronic Computers}, EC-14(3), 326--334.
%
% \bibitem[Efron(1979)]{efron1979}
% Efron, B. (1979). Bootstrap methods: Another look at the jackknife.
% \textit{The Annals of Statistics}, 7(1), 1--26.
%
% \bibitem[Elhage et al.(2021)]{elhage2021}
% Elhage, N., Nanda, N., Olsson, C., Henighan, T., Joseph, N., Mann, B.,
% \ldots\ Olah, C. (2021). \textit{A mathematical framework for transformer
% circuits}. Transformer Circuits Thread.
% \url{https://transformer-circuits.pub/2021/framework/index.html}
%
% \bibitem[Fisher(1936)]{fisher1936}
% Fisher, R. A. (1936). The use of multiple measurements in taxonomic
% problems. \textit{Annals of Eugenics}, 7(2), 179--188.
%
% \bibitem[Spearman(1904)]{spearman1904}
% Spearman, C. (1904). The proof and measurement of association between two
% things. \textit{The American Journal of Psychology}, 15(1), 72--101.
%
% \bibitem[Turner et al.(2023)]{turner2023}
% Turner, A. M., Thiergart, L., Leech, G., Udell, D., Vazquez, J. J., Mini, U.,
% \& MacDiarmid, M. (2023). \textit{Activation addition: Steering language
% models without optimization}. arXiv:2308.10248.
% \url{https://arxiv.org/abs/2308.10248}
%
% \bibitem[Vaswani et al.(2017)]{vaswani2017}
% Vaswani, A., Shazeer, N., Parmar, N., Uszkoreit, J., Jones, L., Gomez, A. N.,
% Kaiser, L., \& Polosukhin, I. (2017). Attention is all you need.
% \textit{Advances in Neural Information Processing Systems}, 30.
%
% \bibitem[Zou et al.(2023)]{zou2023}
% Zou, A., Phan, L., Chen, S., Campbell, J., Guo, P., Ren, R., \ldots\
% Hendrycks, D. (2023). \textit{Representation engineering: A top-down approach
% to AI transparency}. arXiv:2310.01405. \url{https://arxiv.org/abs/2310.01405}
%
% ======================================================================
